% ---------------------------------------------------------------------------
% Společná preambule pro prezentace předmětu Programování (3. ročník)
%
% Vzhled řeší motiv beamerthemeSPSProsek.sty ve stejné složce (převzatý
% z šablony předmětů ZPG/SNT na FS ČVUT a přeznačkovaný na SPŠ na Proseku).
% Tady se jen načte a doplní se metadata a makra předmětu.
%
% Překlad: XeLaTeX (kvůli fontspec a českým znakům)
%   latexmk -xelatex prezentace.tex
% nebo jednoduše  .\nastroje\build_prezentace.ps1
%
% Vkládá se z prezentace v temata/<tema>/prezentace/ takto:
%   % ---------------------------------------------------------------------------
% Společná preambule pro prezentace předmětu Programování (3. ročník)
%
% Vzhled řeší motiv beamerthemeSPSProsek.sty ve stejné složce (převzatý
% z šablony předmětů ZPG/SNT na FS ČVUT a přeznačkovaný na SPŠ na Proseku).
% Tady se jen načte a doplní se metadata a makra předmětu.
%
% Překlad: XeLaTeX (kvůli fontspec a českým znakům)
%   latexmk -xelatex prezentace.tex
% nebo jednoduše  .\nastroje\build_prezentace.ps1
%
% Vkládá se z prezentace v temata/<tema>/prezentace/ takto:
%   % ---------------------------------------------------------------------------
% Společná preambule pro prezentace předmětu Programování (3. ročník)
%
% Vzhled řeší motiv beamerthemeSPSProsek.sty ve stejné složce (převzatý
% z šablony předmětů ZPG/SNT na FS ČVUT a přeznačkovaný na SPŠ na Proseku).
% Tady se jen načte a doplní se metadata a makra předmětu.
%
% Překlad: XeLaTeX (kvůli fontspec a českým znakům)
%   latexmk -xelatex prezentace.tex
% nebo jednoduše  .\nastroje\build_prezentace.ps1
%
% Vkládá se z prezentace v temata/<tema>/prezentace/ takto:
%   % ---------------------------------------------------------------------------
% Společná preambule pro prezentace předmětu Programování (3. ročník)
%
% Vzhled řeší motiv beamerthemeSPSProsek.sty ve stejné složce (převzatý
% z šablony předmětů ZPG/SNT na FS ČVUT a přeznačkovaný na SPŠ na Proseku).
% Tady se jen načte a doplní se metadata a makra předmětu.
%
% Překlad: XeLaTeX (kvůli fontspec a českým znakům)
%   latexmk -xelatex prezentace.tex
% nebo jednoduše  .\nastroje\build_prezentace.ps1
%
% Vkládá se z prezentace v temata/<tema>/prezentace/ takto:
%   \input{../../../spolecne/styl/preambule}
% ---------------------------------------------------------------------------

\usepackage{iftex}

% --- písma a čeština -------------------------------------------------------
\ifXeTeX
  \usepackage{fontspec}
  \usepackage{polyglossia}
  \setmainlanguage{czech}
  \setmainfont{Latin Modern Roman}
  \setsansfont{Latin Modern Sans}
  \setmonofont{Latin Modern Mono}[Scale=MatchLowercase]
\else
  % Záložní varianta, kdyby někdo přeložil pdfLaTeXem.
  \usepackage[T1]{fontenc}
  \usepackage[utf8]{inputenc}
  \usepackage{lmodern}
  \usepackage[czech]{babel}
\fi

\usepackage{graphicx}
\usepackage{booktabs}
\usepackage{xcolor}
\usepackage{tikz}
\usetikzlibrary{arrows.meta, positioning, shapes.geometric}

% --- motiv SPŠ na Proseku --------------------------------------------------
% Cesta ke složce se stylem a logy. Počítá s tím, že prezentace leží
% v temata/<tema>/prezentace/; zkusí se ale i mělčí zanoření, aby šlo
% přeložit i rozpracovaný soubor jinde v repozitáři.
\IfFileExists{../../../spolecne/styl/beamerthemeSPSProsek.sty}
  {\newcommand{\SPSStylPath}{../../../spolecne/styl}}
  {\IfFileExists{../../spolecne/styl/beamerthemeSPSProsek.sty}
    {\newcommand{\SPSStylPath}{../../spolecne/styl}}
    {\IfFileExists{../spolecne/styl/beamerthemeSPSProsek.sty}
      {\newcommand{\SPSStylPath}{../spolecne/styl}}
      {\newcommand{\SPSStylPath}{.}}}}
\usepackage{\SPSStylPath/beamerthemeSPSProsek}

\setbeamertemplate{footline}[SPSProsek]
\beamertemplatenavigationsymbolsempty

% --- velikost písma ve výpisech kódu ---------------------------------------
% Výchozí je \footnotesize. Když se dlouhá ukázka nevejde na slide, zmenši
% ji buď globálně tady, nebo jen v jedné prezentaci stejným řádkem za
% \input{...preambule}:
% \renewcommand{\SPSCodeSize}{\scriptsize}

% --- zkratky pro vkládání kódu ---------------------------------------------
\newcommand{\kod}[1]{\texttt{\small #1}}          % kód uvnitř věty
\newcommand{\kodSoubor}[1]{\lstinputlisting{#1}}  % celá skica ze souboru
\newcommand{\kodSouborRadky}[3]{\lstinputlisting[firstline=#2,lastline=#3]{#1}}

% --- upozornění na slidu ---------------------------------------------------
\newcommand{\pozor}[1]{\textcolor{akcent}{\textbf{#1}}}

% --- metadata předmětu -----------------------------------------------------
\author{Jaroslav Bušek}
\institute{Střední průmyslová škola na Proseku}
% \date{} zůstává prázdné záměrně -- datum a čas posledního překladu se
% zobrazuje automaticky v patě každého slidu (\SPSVersionStamp), takže je
% při promítání hned vidět, jak stará je verze PDF.
\date{}

% ---------------------------------------------------------------------------

\usepackage{iftex}

% --- písma a čeština -------------------------------------------------------
\ifXeTeX
  \usepackage{fontspec}
  \usepackage{polyglossia}
  \setmainlanguage{czech}
  \setmainfont{Latin Modern Roman}
  \setsansfont{Latin Modern Sans}
  \setmonofont{Latin Modern Mono}[Scale=MatchLowercase]
\else
  % Záložní varianta, kdyby někdo přeložil pdfLaTeXem.
  \usepackage[T1]{fontenc}
  \usepackage[utf8]{inputenc}
  \usepackage{lmodern}
  \usepackage[czech]{babel}
\fi

\usepackage{graphicx}
\usepackage{booktabs}
\usepackage{xcolor}
\usepackage{tikz}
\usetikzlibrary{arrows.meta, positioning, shapes.geometric}

% --- motiv SPŠ na Proseku --------------------------------------------------
% Cesta ke složce se stylem a logy. Počítá s tím, že prezentace leží
% v temata/<tema>/prezentace/; zkusí se ale i mělčí zanoření, aby šlo
% přeložit i rozpracovaný soubor jinde v repozitáři.
\IfFileExists{../../../spolecne/styl/beamerthemeSPSProsek.sty}
  {\newcommand{\SPSStylPath}{../../../spolecne/styl}}
  {\IfFileExists{../../spolecne/styl/beamerthemeSPSProsek.sty}
    {\newcommand{\SPSStylPath}{../../spolecne/styl}}
    {\IfFileExists{../spolecne/styl/beamerthemeSPSProsek.sty}
      {\newcommand{\SPSStylPath}{../spolecne/styl}}
      {\newcommand{\SPSStylPath}{.}}}}
\usepackage{\SPSStylPath/beamerthemeSPSProsek}

\setbeamertemplate{footline}[SPSProsek]
\beamertemplatenavigationsymbolsempty

% --- velikost písma ve výpisech kódu ---------------------------------------
% Výchozí je \footnotesize. Když se dlouhá ukázka nevejde na slide, zmenši
% ji buď globálně tady, nebo jen v jedné prezentaci stejným řádkem za
% % ---------------------------------------------------------------------------
% Společná preambule pro prezentace předmětu Programování (3. ročník)
%
% Vzhled řeší motiv beamerthemeSPSProsek.sty ve stejné složce (převzatý
% z šablony předmětů ZPG/SNT na FS ČVUT a přeznačkovaný na SPŠ na Proseku).
% Tady se jen načte a doplní se metadata a makra předmětu.
%
% Překlad: XeLaTeX (kvůli fontspec a českým znakům)
%   latexmk -xelatex prezentace.tex
% nebo jednoduše  .\nastroje\build_prezentace.ps1
%
% Vkládá se z prezentace v temata/<tema>/prezentace/ takto:
%   \input{../../../spolecne/styl/preambule}
% ---------------------------------------------------------------------------

\usepackage{iftex}

% --- písma a čeština -------------------------------------------------------
\ifXeTeX
  \usepackage{fontspec}
  \usepackage{polyglossia}
  \setmainlanguage{czech}
  \setmainfont{Latin Modern Roman}
  \setsansfont{Latin Modern Sans}
  \setmonofont{Latin Modern Mono}[Scale=MatchLowercase]
\else
  % Záložní varianta, kdyby někdo přeložil pdfLaTeXem.
  \usepackage[T1]{fontenc}
  \usepackage[utf8]{inputenc}
  \usepackage{lmodern}
  \usepackage[czech]{babel}
\fi

\usepackage{graphicx}
\usepackage{booktabs}
\usepackage{xcolor}
\usepackage{tikz}
\usetikzlibrary{arrows.meta, positioning, shapes.geometric}

% --- motiv SPŠ na Proseku --------------------------------------------------
% Cesta ke složce se stylem a logy. Počítá s tím, že prezentace leží
% v temata/<tema>/prezentace/; zkusí se ale i mělčí zanoření, aby šlo
% přeložit i rozpracovaný soubor jinde v repozitáři.
\IfFileExists{../../../spolecne/styl/beamerthemeSPSProsek.sty}
  {\newcommand{\SPSStylPath}{../../../spolecne/styl}}
  {\IfFileExists{../../spolecne/styl/beamerthemeSPSProsek.sty}
    {\newcommand{\SPSStylPath}{../../spolecne/styl}}
    {\IfFileExists{../spolecne/styl/beamerthemeSPSProsek.sty}
      {\newcommand{\SPSStylPath}{../spolecne/styl}}
      {\newcommand{\SPSStylPath}{.}}}}
\usepackage{\SPSStylPath/beamerthemeSPSProsek}

\setbeamertemplate{footline}[SPSProsek]
\beamertemplatenavigationsymbolsempty

% --- velikost písma ve výpisech kódu ---------------------------------------
% Výchozí je \footnotesize. Když se dlouhá ukázka nevejde na slide, zmenši
% ji buď globálně tady, nebo jen v jedné prezentaci stejným řádkem za
% \input{...preambule}:
% \renewcommand{\SPSCodeSize}{\scriptsize}

% --- zkratky pro vkládání kódu ---------------------------------------------
\newcommand{\kod}[1]{\texttt{\small #1}}          % kód uvnitř věty
\newcommand{\kodSoubor}[1]{\lstinputlisting{#1}}  % celá skica ze souboru
\newcommand{\kodSouborRadky}[3]{\lstinputlisting[firstline=#2,lastline=#3]{#1}}

% --- upozornění na slidu ---------------------------------------------------
\newcommand{\pozor}[1]{\textcolor{akcent}{\textbf{#1}}}

% --- metadata předmětu -----------------------------------------------------
\author{Jaroslav Bušek}
\institute{Střední průmyslová škola na Proseku}
% \date{} zůstává prázdné záměrně -- datum a čas posledního překladu se
% zobrazuje automaticky v patě každého slidu (\SPSVersionStamp), takže je
% při promítání hned vidět, jak stará je verze PDF.
\date{}
:
% \renewcommand{\SPSCodeSize}{\scriptsize}

% --- zkratky pro vkládání kódu ---------------------------------------------
\newcommand{\kod}[1]{\texttt{\small #1}}          % kód uvnitř věty
\newcommand{\kodSoubor}[1]{\lstinputlisting{#1}}  % celá skica ze souboru
\newcommand{\kodSouborRadky}[3]{\lstinputlisting[firstline=#2,lastline=#3]{#1}}

% --- upozornění na slidu ---------------------------------------------------
\newcommand{\pozor}[1]{\textcolor{akcent}{\textbf{#1}}}

% --- metadata předmětu -----------------------------------------------------
\author{Jaroslav Bušek}
\institute{Střední průmyslová škola na Proseku}
% \date{} zůstává prázdné záměrně -- datum a čas posledního překladu se
% zobrazuje automaticky v patě každého slidu (\SPSVersionStamp), takže je
% při promítání hned vidět, jak stará je verze PDF.
\date{}

% ---------------------------------------------------------------------------

\usepackage{iftex}

% --- písma a čeština -------------------------------------------------------
\ifXeTeX
  \usepackage{fontspec}
  \usepackage{polyglossia}
  \setmainlanguage{czech}
  \setmainfont{Latin Modern Roman}
  \setsansfont{Latin Modern Sans}
  \setmonofont{Latin Modern Mono}[Scale=MatchLowercase]
\else
  % Záložní varianta, kdyby někdo přeložil pdfLaTeXem.
  \usepackage[T1]{fontenc}
  \usepackage[utf8]{inputenc}
  \usepackage{lmodern}
  \usepackage[czech]{babel}
\fi

\usepackage{graphicx}
\usepackage{booktabs}
\usepackage{xcolor}
\usepackage{tikz}
\usetikzlibrary{arrows.meta, positioning, shapes.geometric}

% --- motiv SPŠ na Proseku --------------------------------------------------
% Cesta ke složce se stylem a logy. Počítá s tím, že prezentace leží
% v temata/<tema>/prezentace/; zkusí se ale i mělčí zanoření, aby šlo
% přeložit i rozpracovaný soubor jinde v repozitáři.
\IfFileExists{../../../spolecne/styl/beamerthemeSPSProsek.sty}
  {\newcommand{\SPSStylPath}{../../../spolecne/styl}}
  {\IfFileExists{../../spolecne/styl/beamerthemeSPSProsek.sty}
    {\newcommand{\SPSStylPath}{../../spolecne/styl}}
    {\IfFileExists{../spolecne/styl/beamerthemeSPSProsek.sty}
      {\newcommand{\SPSStylPath}{../spolecne/styl}}
      {\newcommand{\SPSStylPath}{.}}}}
\usepackage{\SPSStylPath/beamerthemeSPSProsek}

\setbeamertemplate{footline}[SPSProsek]
\beamertemplatenavigationsymbolsempty

% --- velikost písma ve výpisech kódu ---------------------------------------
% Výchozí je \footnotesize. Když se dlouhá ukázka nevejde na slide, zmenši
% ji buď globálně tady, nebo jen v jedné prezentaci stejným řádkem za
% % ---------------------------------------------------------------------------
% Společná preambule pro prezentace předmětu Programování (3. ročník)
%
% Vzhled řeší motiv beamerthemeSPSProsek.sty ve stejné složce (převzatý
% z šablony předmětů ZPG/SNT na FS ČVUT a přeznačkovaný na SPŠ na Proseku).
% Tady se jen načte a doplní se metadata a makra předmětu.
%
% Překlad: XeLaTeX (kvůli fontspec a českým znakům)
%   latexmk -xelatex prezentace.tex
% nebo jednoduše  .\nastroje\build_prezentace.ps1
%
% Vkládá se z prezentace v temata/<tema>/prezentace/ takto:
%   % ---------------------------------------------------------------------------
% Společná preambule pro prezentace předmětu Programování (3. ročník)
%
% Vzhled řeší motiv beamerthemeSPSProsek.sty ve stejné složce (převzatý
% z šablony předmětů ZPG/SNT na FS ČVUT a přeznačkovaný na SPŠ na Proseku).
% Tady se jen načte a doplní se metadata a makra předmětu.
%
% Překlad: XeLaTeX (kvůli fontspec a českým znakům)
%   latexmk -xelatex prezentace.tex
% nebo jednoduše  .\nastroje\build_prezentace.ps1
%
% Vkládá se z prezentace v temata/<tema>/prezentace/ takto:
%   \input{../../../spolecne/styl/preambule}
% ---------------------------------------------------------------------------

\usepackage{iftex}

% --- písma a čeština -------------------------------------------------------
\ifXeTeX
  \usepackage{fontspec}
  \usepackage{polyglossia}
  \setmainlanguage{czech}
  \setmainfont{Latin Modern Roman}
  \setsansfont{Latin Modern Sans}
  \setmonofont{Latin Modern Mono}[Scale=MatchLowercase]
\else
  % Záložní varianta, kdyby někdo přeložil pdfLaTeXem.
  \usepackage[T1]{fontenc}
  \usepackage[utf8]{inputenc}
  \usepackage{lmodern}
  \usepackage[czech]{babel}
\fi

\usepackage{graphicx}
\usepackage{booktabs}
\usepackage{xcolor}
\usepackage{tikz}
\usetikzlibrary{arrows.meta, positioning, shapes.geometric}

% --- motiv SPŠ na Proseku --------------------------------------------------
% Cesta ke složce se stylem a logy. Počítá s tím, že prezentace leží
% v temata/<tema>/prezentace/; zkusí se ale i mělčí zanoření, aby šlo
% přeložit i rozpracovaný soubor jinde v repozitáři.
\IfFileExists{../../../spolecne/styl/beamerthemeSPSProsek.sty}
  {\newcommand{\SPSStylPath}{../../../spolecne/styl}}
  {\IfFileExists{../../spolecne/styl/beamerthemeSPSProsek.sty}
    {\newcommand{\SPSStylPath}{../../spolecne/styl}}
    {\IfFileExists{../spolecne/styl/beamerthemeSPSProsek.sty}
      {\newcommand{\SPSStylPath}{../spolecne/styl}}
      {\newcommand{\SPSStylPath}{.}}}}
\usepackage{\SPSStylPath/beamerthemeSPSProsek}

\setbeamertemplate{footline}[SPSProsek]
\beamertemplatenavigationsymbolsempty

% --- velikost písma ve výpisech kódu ---------------------------------------
% Výchozí je \footnotesize. Když se dlouhá ukázka nevejde na slide, zmenši
% ji buď globálně tady, nebo jen v jedné prezentaci stejným řádkem za
% \input{...preambule}:
% \renewcommand{\SPSCodeSize}{\scriptsize}

% --- zkratky pro vkládání kódu ---------------------------------------------
\newcommand{\kod}[1]{\texttt{\small #1}}          % kód uvnitř věty
\newcommand{\kodSoubor}[1]{\lstinputlisting{#1}}  % celá skica ze souboru
\newcommand{\kodSouborRadky}[3]{\lstinputlisting[firstline=#2,lastline=#3]{#1}}

% --- upozornění na slidu ---------------------------------------------------
\newcommand{\pozor}[1]{\textcolor{akcent}{\textbf{#1}}}

% --- metadata předmětu -----------------------------------------------------
\author{Jaroslav Bušek}
\institute{Střední průmyslová škola na Proseku}
% \date{} zůstává prázdné záměrně -- datum a čas posledního překladu se
% zobrazuje automaticky v patě každého slidu (\SPSVersionStamp), takže je
% při promítání hned vidět, jak stará je verze PDF.
\date{}

% ---------------------------------------------------------------------------

\usepackage{iftex}

% --- písma a čeština -------------------------------------------------------
\ifXeTeX
  \usepackage{fontspec}
  \usepackage{polyglossia}
  \setmainlanguage{czech}
  \setmainfont{Latin Modern Roman}
  \setsansfont{Latin Modern Sans}
  \setmonofont{Latin Modern Mono}[Scale=MatchLowercase]
\else
  % Záložní varianta, kdyby někdo přeložil pdfLaTeXem.
  \usepackage[T1]{fontenc}
  \usepackage[utf8]{inputenc}
  \usepackage{lmodern}
  \usepackage[czech]{babel}
\fi

\usepackage{graphicx}
\usepackage{booktabs}
\usepackage{xcolor}
\usepackage{tikz}
\usetikzlibrary{arrows.meta, positioning, shapes.geometric}

% --- motiv SPŠ na Proseku --------------------------------------------------
% Cesta ke složce se stylem a logy. Počítá s tím, že prezentace leží
% v temata/<tema>/prezentace/; zkusí se ale i mělčí zanoření, aby šlo
% přeložit i rozpracovaný soubor jinde v repozitáři.
\IfFileExists{../../../spolecne/styl/beamerthemeSPSProsek.sty}
  {\newcommand{\SPSStylPath}{../../../spolecne/styl}}
  {\IfFileExists{../../spolecne/styl/beamerthemeSPSProsek.sty}
    {\newcommand{\SPSStylPath}{../../spolecne/styl}}
    {\IfFileExists{../spolecne/styl/beamerthemeSPSProsek.sty}
      {\newcommand{\SPSStylPath}{../spolecne/styl}}
      {\newcommand{\SPSStylPath}{.}}}}
\usepackage{\SPSStylPath/beamerthemeSPSProsek}

\setbeamertemplate{footline}[SPSProsek]
\beamertemplatenavigationsymbolsempty

% --- velikost písma ve výpisech kódu ---------------------------------------
% Výchozí je \footnotesize. Když se dlouhá ukázka nevejde na slide, zmenši
% ji buď globálně tady, nebo jen v jedné prezentaci stejným řádkem za
% % ---------------------------------------------------------------------------
% Společná preambule pro prezentace předmětu Programování (3. ročník)
%
% Vzhled řeší motiv beamerthemeSPSProsek.sty ve stejné složce (převzatý
% z šablony předmětů ZPG/SNT na FS ČVUT a přeznačkovaný na SPŠ na Proseku).
% Tady se jen načte a doplní se metadata a makra předmětu.
%
% Překlad: XeLaTeX (kvůli fontspec a českým znakům)
%   latexmk -xelatex prezentace.tex
% nebo jednoduše  .\nastroje\build_prezentace.ps1
%
% Vkládá se z prezentace v temata/<tema>/prezentace/ takto:
%   \input{../../../spolecne/styl/preambule}
% ---------------------------------------------------------------------------

\usepackage{iftex}

% --- písma a čeština -------------------------------------------------------
\ifXeTeX
  \usepackage{fontspec}
  \usepackage{polyglossia}
  \setmainlanguage{czech}
  \setmainfont{Latin Modern Roman}
  \setsansfont{Latin Modern Sans}
  \setmonofont{Latin Modern Mono}[Scale=MatchLowercase]
\else
  % Záložní varianta, kdyby někdo přeložil pdfLaTeXem.
  \usepackage[T1]{fontenc}
  \usepackage[utf8]{inputenc}
  \usepackage{lmodern}
  \usepackage[czech]{babel}
\fi

\usepackage{graphicx}
\usepackage{booktabs}
\usepackage{xcolor}
\usepackage{tikz}
\usetikzlibrary{arrows.meta, positioning, shapes.geometric}

% --- motiv SPŠ na Proseku --------------------------------------------------
% Cesta ke složce se stylem a logy. Počítá s tím, že prezentace leží
% v temata/<tema>/prezentace/; zkusí se ale i mělčí zanoření, aby šlo
% přeložit i rozpracovaný soubor jinde v repozitáři.
\IfFileExists{../../../spolecne/styl/beamerthemeSPSProsek.sty}
  {\newcommand{\SPSStylPath}{../../../spolecne/styl}}
  {\IfFileExists{../../spolecne/styl/beamerthemeSPSProsek.sty}
    {\newcommand{\SPSStylPath}{../../spolecne/styl}}
    {\IfFileExists{../spolecne/styl/beamerthemeSPSProsek.sty}
      {\newcommand{\SPSStylPath}{../spolecne/styl}}
      {\newcommand{\SPSStylPath}{.}}}}
\usepackage{\SPSStylPath/beamerthemeSPSProsek}

\setbeamertemplate{footline}[SPSProsek]
\beamertemplatenavigationsymbolsempty

% --- velikost písma ve výpisech kódu ---------------------------------------
% Výchozí je \footnotesize. Když se dlouhá ukázka nevejde na slide, zmenši
% ji buď globálně tady, nebo jen v jedné prezentaci stejným řádkem za
% \input{...preambule}:
% \renewcommand{\SPSCodeSize}{\scriptsize}

% --- zkratky pro vkládání kódu ---------------------------------------------
\newcommand{\kod}[1]{\texttt{\small #1}}          % kód uvnitř věty
\newcommand{\kodSoubor}[1]{\lstinputlisting{#1}}  % celá skica ze souboru
\newcommand{\kodSouborRadky}[3]{\lstinputlisting[firstline=#2,lastline=#3]{#1}}

% --- upozornění na slidu ---------------------------------------------------
\newcommand{\pozor}[1]{\textcolor{akcent}{\textbf{#1}}}

% --- metadata předmětu -----------------------------------------------------
\author{Jaroslav Bušek}
\institute{Střední průmyslová škola na Proseku}
% \date{} zůstává prázdné záměrně -- datum a čas posledního překladu se
% zobrazuje automaticky v patě každého slidu (\SPSVersionStamp), takže je
% při promítání hned vidět, jak stará je verze PDF.
\date{}
:
% \renewcommand{\SPSCodeSize}{\scriptsize}

% --- zkratky pro vkládání kódu ---------------------------------------------
\newcommand{\kod}[1]{\texttt{\small #1}}          % kód uvnitř věty
\newcommand{\kodSoubor}[1]{\lstinputlisting{#1}}  % celá skica ze souboru
\newcommand{\kodSouborRadky}[3]{\lstinputlisting[firstline=#2,lastline=#3]{#1}}

% --- upozornění na slidu ---------------------------------------------------
\newcommand{\pozor}[1]{\textcolor{akcent}{\textbf{#1}}}

% --- metadata předmětu -----------------------------------------------------
\author{Jaroslav Bušek}
\institute{Střední průmyslová škola na Proseku}
% \date{} zůstává prázdné záměrně -- datum a čas posledního překladu se
% zobrazuje automaticky v patě každého slidu (\SPSVersionStamp), takže je
% při promítání hned vidět, jak stará je verze PDF.
\date{}
:
% \renewcommand{\SPSCodeSize}{\scriptsize}

% --- zkratky pro vkládání kódu ---------------------------------------------
\newcommand{\kod}[1]{\texttt{\small #1}}          % kód uvnitř věty
\newcommand{\kodSoubor}[1]{\lstinputlisting{#1}}  % celá skica ze souboru
\newcommand{\kodSouborRadky}[3]{\lstinputlisting[firstline=#2,lastline=#3]{#1}}

% --- upozornění na slidu ---------------------------------------------------
\newcommand{\pozor}[1]{\textcolor{akcent}{\textbf{#1}}}

% --- metadata předmětu -----------------------------------------------------
\author{Jaroslav Bušek}
\institute{Střední průmyslová škola na Proseku}
% \date{} zůstává prázdné záměrně -- datum a čas posledního překladu se
% zobrazuje automaticky v patě každého slidu (\SPSVersionStamp), takže je
% při promítání hned vidět, jak stará je verze PDF.
\date{}

% ---------------------------------------------------------------------------

\usepackage{iftex}

% --- písma a čeština -------------------------------------------------------
\ifXeTeX
  \usepackage{fontspec}
  \usepackage{polyglossia}
  \setmainlanguage{czech}
  \setmainfont{Latin Modern Roman}
  \setsansfont{Latin Modern Sans}
  \setmonofont{Latin Modern Mono}[Scale=MatchLowercase]
\else
  % Záložní varianta, kdyby někdo přeložil pdfLaTeXem.
  \usepackage[T1]{fontenc}
  \usepackage[utf8]{inputenc}
  \usepackage{lmodern}
  \usepackage[czech]{babel}
\fi

\usepackage{graphicx}
\usepackage{booktabs}
\usepackage{xcolor}
\usepackage{tikz}
\usetikzlibrary{arrows.meta, positioning, shapes.geometric}

% --- motiv SPŠ na Proseku --------------------------------------------------
% Cesta ke složce se stylem a logy. Počítá s tím, že prezentace leží
% v temata/<tema>/prezentace/; zkusí se ale i mělčí zanoření, aby šlo
% přeložit i rozpracovaný soubor jinde v repozitáři.
\IfFileExists{../../../spolecne/styl/beamerthemeSPSProsek.sty}
  {\newcommand{\SPSStylPath}{../../../spolecne/styl}}
  {\IfFileExists{../../spolecne/styl/beamerthemeSPSProsek.sty}
    {\newcommand{\SPSStylPath}{../../spolecne/styl}}
    {\IfFileExists{../spolecne/styl/beamerthemeSPSProsek.sty}
      {\newcommand{\SPSStylPath}{../spolecne/styl}}
      {\newcommand{\SPSStylPath}{.}}}}
\usepackage{\SPSStylPath/beamerthemeSPSProsek}

\setbeamertemplate{footline}[SPSProsek]
\beamertemplatenavigationsymbolsempty

% --- velikost písma ve výpisech kódu ---------------------------------------
% Výchozí je \footnotesize. Když se dlouhá ukázka nevejde na slide, zmenši
% ji buď globálně tady, nebo jen v jedné prezentaci stejným řádkem za
% % ---------------------------------------------------------------------------
% Společná preambule pro prezentace předmětu Programování (3. ročník)
%
% Vzhled řeší motiv beamerthemeSPSProsek.sty ve stejné složce (převzatý
% z šablony předmětů ZPG/SNT na FS ČVUT a přeznačkovaný na SPŠ na Proseku).
% Tady se jen načte a doplní se metadata a makra předmětu.
%
% Překlad: XeLaTeX (kvůli fontspec a českým znakům)
%   latexmk -xelatex prezentace.tex
% nebo jednoduše  .\nastroje\build_prezentace.ps1
%
% Vkládá se z prezentace v temata/<tema>/prezentace/ takto:
%   % ---------------------------------------------------------------------------
% Společná preambule pro prezentace předmětu Programování (3. ročník)
%
% Vzhled řeší motiv beamerthemeSPSProsek.sty ve stejné složce (převzatý
% z šablony předmětů ZPG/SNT na FS ČVUT a přeznačkovaný na SPŠ na Proseku).
% Tady se jen načte a doplní se metadata a makra předmětu.
%
% Překlad: XeLaTeX (kvůli fontspec a českým znakům)
%   latexmk -xelatex prezentace.tex
% nebo jednoduše  .\nastroje\build_prezentace.ps1
%
% Vkládá se z prezentace v temata/<tema>/prezentace/ takto:
%   % ---------------------------------------------------------------------------
% Společná preambule pro prezentace předmětu Programování (3. ročník)
%
% Vzhled řeší motiv beamerthemeSPSProsek.sty ve stejné složce (převzatý
% z šablony předmětů ZPG/SNT na FS ČVUT a přeznačkovaný na SPŠ na Proseku).
% Tady se jen načte a doplní se metadata a makra předmětu.
%
% Překlad: XeLaTeX (kvůli fontspec a českým znakům)
%   latexmk -xelatex prezentace.tex
% nebo jednoduše  .\nastroje\build_prezentace.ps1
%
% Vkládá se z prezentace v temata/<tema>/prezentace/ takto:
%   \input{../../../spolecne/styl/preambule}
% ---------------------------------------------------------------------------

\usepackage{iftex}

% --- písma a čeština -------------------------------------------------------
\ifXeTeX
  \usepackage{fontspec}
  \usepackage{polyglossia}
  \setmainlanguage{czech}
  \setmainfont{Latin Modern Roman}
  \setsansfont{Latin Modern Sans}
  \setmonofont{Latin Modern Mono}[Scale=MatchLowercase]
\else
  % Záložní varianta, kdyby někdo přeložil pdfLaTeXem.
  \usepackage[T1]{fontenc}
  \usepackage[utf8]{inputenc}
  \usepackage{lmodern}
  \usepackage[czech]{babel}
\fi

\usepackage{graphicx}
\usepackage{booktabs}
\usepackage{xcolor}
\usepackage{tikz}
\usetikzlibrary{arrows.meta, positioning, shapes.geometric}

% --- motiv SPŠ na Proseku --------------------------------------------------
% Cesta ke složce se stylem a logy. Počítá s tím, že prezentace leží
% v temata/<tema>/prezentace/; zkusí se ale i mělčí zanoření, aby šlo
% přeložit i rozpracovaný soubor jinde v repozitáři.
\IfFileExists{../../../spolecne/styl/beamerthemeSPSProsek.sty}
  {\newcommand{\SPSStylPath}{../../../spolecne/styl}}
  {\IfFileExists{../../spolecne/styl/beamerthemeSPSProsek.sty}
    {\newcommand{\SPSStylPath}{../../spolecne/styl}}
    {\IfFileExists{../spolecne/styl/beamerthemeSPSProsek.sty}
      {\newcommand{\SPSStylPath}{../spolecne/styl}}
      {\newcommand{\SPSStylPath}{.}}}}
\usepackage{\SPSStylPath/beamerthemeSPSProsek}

\setbeamertemplate{footline}[SPSProsek]
\beamertemplatenavigationsymbolsempty

% --- velikost písma ve výpisech kódu ---------------------------------------
% Výchozí je \footnotesize. Když se dlouhá ukázka nevejde na slide, zmenši
% ji buď globálně tady, nebo jen v jedné prezentaci stejným řádkem za
% \input{...preambule}:
% \renewcommand{\SPSCodeSize}{\scriptsize}

% --- zkratky pro vkládání kódu ---------------------------------------------
\newcommand{\kod}[1]{\texttt{\small #1}}          % kód uvnitř věty
\newcommand{\kodSoubor}[1]{\lstinputlisting{#1}}  % celá skica ze souboru
\newcommand{\kodSouborRadky}[3]{\lstinputlisting[firstline=#2,lastline=#3]{#1}}

% --- upozornění na slidu ---------------------------------------------------
\newcommand{\pozor}[1]{\textcolor{akcent}{\textbf{#1}}}

% --- metadata předmětu -----------------------------------------------------
\author{Jaroslav Bušek}
\institute{Střední průmyslová škola na Proseku}
% \date{} zůstává prázdné záměrně -- datum a čas posledního překladu se
% zobrazuje automaticky v patě každého slidu (\SPSVersionStamp), takže je
% při promítání hned vidět, jak stará je verze PDF.
\date{}

% ---------------------------------------------------------------------------

\usepackage{iftex}

% --- písma a čeština -------------------------------------------------------
\ifXeTeX
  \usepackage{fontspec}
  \usepackage{polyglossia}
  \setmainlanguage{czech}
  \setmainfont{Latin Modern Roman}
  \setsansfont{Latin Modern Sans}
  \setmonofont{Latin Modern Mono}[Scale=MatchLowercase]
\else
  % Záložní varianta, kdyby někdo přeložil pdfLaTeXem.
  \usepackage[T1]{fontenc}
  \usepackage[utf8]{inputenc}
  \usepackage{lmodern}
  \usepackage[czech]{babel}
\fi

\usepackage{graphicx}
\usepackage{booktabs}
\usepackage{xcolor}
\usepackage{tikz}
\usetikzlibrary{arrows.meta, positioning, shapes.geometric}

% --- motiv SPŠ na Proseku --------------------------------------------------
% Cesta ke složce se stylem a logy. Počítá s tím, že prezentace leží
% v temata/<tema>/prezentace/; zkusí se ale i mělčí zanoření, aby šlo
% přeložit i rozpracovaný soubor jinde v repozitáři.
\IfFileExists{../../../spolecne/styl/beamerthemeSPSProsek.sty}
  {\newcommand{\SPSStylPath}{../../../spolecne/styl}}
  {\IfFileExists{../../spolecne/styl/beamerthemeSPSProsek.sty}
    {\newcommand{\SPSStylPath}{../../spolecne/styl}}
    {\IfFileExists{../spolecne/styl/beamerthemeSPSProsek.sty}
      {\newcommand{\SPSStylPath}{../spolecne/styl}}
      {\newcommand{\SPSStylPath}{.}}}}
\usepackage{\SPSStylPath/beamerthemeSPSProsek}

\setbeamertemplate{footline}[SPSProsek]
\beamertemplatenavigationsymbolsempty

% --- velikost písma ve výpisech kódu ---------------------------------------
% Výchozí je \footnotesize. Když se dlouhá ukázka nevejde na slide, zmenši
% ji buď globálně tady, nebo jen v jedné prezentaci stejným řádkem za
% % ---------------------------------------------------------------------------
% Společná preambule pro prezentace předmětu Programování (3. ročník)
%
% Vzhled řeší motiv beamerthemeSPSProsek.sty ve stejné složce (převzatý
% z šablony předmětů ZPG/SNT na FS ČVUT a přeznačkovaný na SPŠ na Proseku).
% Tady se jen načte a doplní se metadata a makra předmětu.
%
% Překlad: XeLaTeX (kvůli fontspec a českým znakům)
%   latexmk -xelatex prezentace.tex
% nebo jednoduše  .\nastroje\build_prezentace.ps1
%
% Vkládá se z prezentace v temata/<tema>/prezentace/ takto:
%   \input{../../../spolecne/styl/preambule}
% ---------------------------------------------------------------------------

\usepackage{iftex}

% --- písma a čeština -------------------------------------------------------
\ifXeTeX
  \usepackage{fontspec}
  \usepackage{polyglossia}
  \setmainlanguage{czech}
  \setmainfont{Latin Modern Roman}
  \setsansfont{Latin Modern Sans}
  \setmonofont{Latin Modern Mono}[Scale=MatchLowercase]
\else
  % Záložní varianta, kdyby někdo přeložil pdfLaTeXem.
  \usepackage[T1]{fontenc}
  \usepackage[utf8]{inputenc}
  \usepackage{lmodern}
  \usepackage[czech]{babel}
\fi

\usepackage{graphicx}
\usepackage{booktabs}
\usepackage{xcolor}
\usepackage{tikz}
\usetikzlibrary{arrows.meta, positioning, shapes.geometric}

% --- motiv SPŠ na Proseku --------------------------------------------------
% Cesta ke složce se stylem a logy. Počítá s tím, že prezentace leží
% v temata/<tema>/prezentace/; zkusí se ale i mělčí zanoření, aby šlo
% přeložit i rozpracovaný soubor jinde v repozitáři.
\IfFileExists{../../../spolecne/styl/beamerthemeSPSProsek.sty}
  {\newcommand{\SPSStylPath}{../../../spolecne/styl}}
  {\IfFileExists{../../spolecne/styl/beamerthemeSPSProsek.sty}
    {\newcommand{\SPSStylPath}{../../spolecne/styl}}
    {\IfFileExists{../spolecne/styl/beamerthemeSPSProsek.sty}
      {\newcommand{\SPSStylPath}{../spolecne/styl}}
      {\newcommand{\SPSStylPath}{.}}}}
\usepackage{\SPSStylPath/beamerthemeSPSProsek}

\setbeamertemplate{footline}[SPSProsek]
\beamertemplatenavigationsymbolsempty

% --- velikost písma ve výpisech kódu ---------------------------------------
% Výchozí je \footnotesize. Když se dlouhá ukázka nevejde na slide, zmenši
% ji buď globálně tady, nebo jen v jedné prezentaci stejným řádkem za
% \input{...preambule}:
% \renewcommand{\SPSCodeSize}{\scriptsize}

% --- zkratky pro vkládání kódu ---------------------------------------------
\newcommand{\kod}[1]{\texttt{\small #1}}          % kód uvnitř věty
\newcommand{\kodSoubor}[1]{\lstinputlisting{#1}}  % celá skica ze souboru
\newcommand{\kodSouborRadky}[3]{\lstinputlisting[firstline=#2,lastline=#3]{#1}}

% --- upozornění na slidu ---------------------------------------------------
\newcommand{\pozor}[1]{\textcolor{akcent}{\textbf{#1}}}

% --- metadata předmětu -----------------------------------------------------
\author{Jaroslav Bušek}
\institute{Střední průmyslová škola na Proseku}
% \date{} zůstává prázdné záměrně -- datum a čas posledního překladu se
% zobrazuje automaticky v patě každého slidu (\SPSVersionStamp), takže je
% při promítání hned vidět, jak stará je verze PDF.
\date{}
:
% \renewcommand{\SPSCodeSize}{\scriptsize}

% --- zkratky pro vkládání kódu ---------------------------------------------
\newcommand{\kod}[1]{\texttt{\small #1}}          % kód uvnitř věty
\newcommand{\kodSoubor}[1]{\lstinputlisting{#1}}  % celá skica ze souboru
\newcommand{\kodSouborRadky}[3]{\lstinputlisting[firstline=#2,lastline=#3]{#1}}

% --- upozornění na slidu ---------------------------------------------------
\newcommand{\pozor}[1]{\textcolor{akcent}{\textbf{#1}}}

% --- metadata předmětu -----------------------------------------------------
\author{Jaroslav Bušek}
\institute{Střední průmyslová škola na Proseku}
% \date{} zůstává prázdné záměrně -- datum a čas posledního překladu se
% zobrazuje automaticky v patě každého slidu (\SPSVersionStamp), takže je
% při promítání hned vidět, jak stará je verze PDF.
\date{}

% ---------------------------------------------------------------------------

\usepackage{iftex}

% --- písma a čeština -------------------------------------------------------
\ifXeTeX
  \usepackage{fontspec}
  \usepackage{polyglossia}
  \setmainlanguage{czech}
  \setmainfont{Latin Modern Roman}
  \setsansfont{Latin Modern Sans}
  \setmonofont{Latin Modern Mono}[Scale=MatchLowercase]
\else
  % Záložní varianta, kdyby někdo přeložil pdfLaTeXem.
  \usepackage[T1]{fontenc}
  \usepackage[utf8]{inputenc}
  \usepackage{lmodern}
  \usepackage[czech]{babel}
\fi

\usepackage{graphicx}
\usepackage{booktabs}
\usepackage{xcolor}
\usepackage{tikz}
\usetikzlibrary{arrows.meta, positioning, shapes.geometric}

% --- motiv SPŠ na Proseku --------------------------------------------------
% Cesta ke složce se stylem a logy. Počítá s tím, že prezentace leží
% v temata/<tema>/prezentace/; zkusí se ale i mělčí zanoření, aby šlo
% přeložit i rozpracovaný soubor jinde v repozitáři.
\IfFileExists{../../../spolecne/styl/beamerthemeSPSProsek.sty}
  {\newcommand{\SPSStylPath}{../../../spolecne/styl}}
  {\IfFileExists{../../spolecne/styl/beamerthemeSPSProsek.sty}
    {\newcommand{\SPSStylPath}{../../spolecne/styl}}
    {\IfFileExists{../spolecne/styl/beamerthemeSPSProsek.sty}
      {\newcommand{\SPSStylPath}{../spolecne/styl}}
      {\newcommand{\SPSStylPath}{.}}}}
\usepackage{\SPSStylPath/beamerthemeSPSProsek}

\setbeamertemplate{footline}[SPSProsek]
\beamertemplatenavigationsymbolsempty

% --- velikost písma ve výpisech kódu ---------------------------------------
% Výchozí je \footnotesize. Když se dlouhá ukázka nevejde na slide, zmenši
% ji buď globálně tady, nebo jen v jedné prezentaci stejným řádkem za
% % ---------------------------------------------------------------------------
% Společná preambule pro prezentace předmětu Programování (3. ročník)
%
% Vzhled řeší motiv beamerthemeSPSProsek.sty ve stejné složce (převzatý
% z šablony předmětů ZPG/SNT na FS ČVUT a přeznačkovaný na SPŠ na Proseku).
% Tady se jen načte a doplní se metadata a makra předmětu.
%
% Překlad: XeLaTeX (kvůli fontspec a českým znakům)
%   latexmk -xelatex prezentace.tex
% nebo jednoduše  .\nastroje\build_prezentace.ps1
%
% Vkládá se z prezentace v temata/<tema>/prezentace/ takto:
%   % ---------------------------------------------------------------------------
% Společná preambule pro prezentace předmětu Programování (3. ročník)
%
% Vzhled řeší motiv beamerthemeSPSProsek.sty ve stejné složce (převzatý
% z šablony předmětů ZPG/SNT na FS ČVUT a přeznačkovaný na SPŠ na Proseku).
% Tady se jen načte a doplní se metadata a makra předmětu.
%
% Překlad: XeLaTeX (kvůli fontspec a českým znakům)
%   latexmk -xelatex prezentace.tex
% nebo jednoduše  .\nastroje\build_prezentace.ps1
%
% Vkládá se z prezentace v temata/<tema>/prezentace/ takto:
%   \input{../../../spolecne/styl/preambule}
% ---------------------------------------------------------------------------

\usepackage{iftex}

% --- písma a čeština -------------------------------------------------------
\ifXeTeX
  \usepackage{fontspec}
  \usepackage{polyglossia}
  \setmainlanguage{czech}
  \setmainfont{Latin Modern Roman}
  \setsansfont{Latin Modern Sans}
  \setmonofont{Latin Modern Mono}[Scale=MatchLowercase]
\else
  % Záložní varianta, kdyby někdo přeložil pdfLaTeXem.
  \usepackage[T1]{fontenc}
  \usepackage[utf8]{inputenc}
  \usepackage{lmodern}
  \usepackage[czech]{babel}
\fi

\usepackage{graphicx}
\usepackage{booktabs}
\usepackage{xcolor}
\usepackage{tikz}
\usetikzlibrary{arrows.meta, positioning, shapes.geometric}

% --- motiv SPŠ na Proseku --------------------------------------------------
% Cesta ke složce se stylem a logy. Počítá s tím, že prezentace leží
% v temata/<tema>/prezentace/; zkusí se ale i mělčí zanoření, aby šlo
% přeložit i rozpracovaný soubor jinde v repozitáři.
\IfFileExists{../../../spolecne/styl/beamerthemeSPSProsek.sty}
  {\newcommand{\SPSStylPath}{../../../spolecne/styl}}
  {\IfFileExists{../../spolecne/styl/beamerthemeSPSProsek.sty}
    {\newcommand{\SPSStylPath}{../../spolecne/styl}}
    {\IfFileExists{../spolecne/styl/beamerthemeSPSProsek.sty}
      {\newcommand{\SPSStylPath}{../spolecne/styl}}
      {\newcommand{\SPSStylPath}{.}}}}
\usepackage{\SPSStylPath/beamerthemeSPSProsek}

\setbeamertemplate{footline}[SPSProsek]
\beamertemplatenavigationsymbolsempty

% --- velikost písma ve výpisech kódu ---------------------------------------
% Výchozí je \footnotesize. Když se dlouhá ukázka nevejde na slide, zmenši
% ji buď globálně tady, nebo jen v jedné prezentaci stejným řádkem za
% \input{...preambule}:
% \renewcommand{\SPSCodeSize}{\scriptsize}

% --- zkratky pro vkládání kódu ---------------------------------------------
\newcommand{\kod}[1]{\texttt{\small #1}}          % kód uvnitř věty
\newcommand{\kodSoubor}[1]{\lstinputlisting{#1}}  % celá skica ze souboru
\newcommand{\kodSouborRadky}[3]{\lstinputlisting[firstline=#2,lastline=#3]{#1}}

% --- upozornění na slidu ---------------------------------------------------
\newcommand{\pozor}[1]{\textcolor{akcent}{\textbf{#1}}}

% --- metadata předmětu -----------------------------------------------------
\author{Jaroslav Bušek}
\institute{Střední průmyslová škola na Proseku}
% \date{} zůstává prázdné záměrně -- datum a čas posledního překladu se
% zobrazuje automaticky v patě každého slidu (\SPSVersionStamp), takže je
% při promítání hned vidět, jak stará je verze PDF.
\date{}

% ---------------------------------------------------------------------------

\usepackage{iftex}

% --- písma a čeština -------------------------------------------------------
\ifXeTeX
  \usepackage{fontspec}
  \usepackage{polyglossia}
  \setmainlanguage{czech}
  \setmainfont{Latin Modern Roman}
  \setsansfont{Latin Modern Sans}
  \setmonofont{Latin Modern Mono}[Scale=MatchLowercase]
\else
  % Záložní varianta, kdyby někdo přeložil pdfLaTeXem.
  \usepackage[T1]{fontenc}
  \usepackage[utf8]{inputenc}
  \usepackage{lmodern}
  \usepackage[czech]{babel}
\fi

\usepackage{graphicx}
\usepackage{booktabs}
\usepackage{xcolor}
\usepackage{tikz}
\usetikzlibrary{arrows.meta, positioning, shapes.geometric}

% --- motiv SPŠ na Proseku --------------------------------------------------
% Cesta ke složce se stylem a logy. Počítá s tím, že prezentace leží
% v temata/<tema>/prezentace/; zkusí se ale i mělčí zanoření, aby šlo
% přeložit i rozpracovaný soubor jinde v repozitáři.
\IfFileExists{../../../spolecne/styl/beamerthemeSPSProsek.sty}
  {\newcommand{\SPSStylPath}{../../../spolecne/styl}}
  {\IfFileExists{../../spolecne/styl/beamerthemeSPSProsek.sty}
    {\newcommand{\SPSStylPath}{../../spolecne/styl}}
    {\IfFileExists{../spolecne/styl/beamerthemeSPSProsek.sty}
      {\newcommand{\SPSStylPath}{../spolecne/styl}}
      {\newcommand{\SPSStylPath}{.}}}}
\usepackage{\SPSStylPath/beamerthemeSPSProsek}

\setbeamertemplate{footline}[SPSProsek]
\beamertemplatenavigationsymbolsempty

% --- velikost písma ve výpisech kódu ---------------------------------------
% Výchozí je \footnotesize. Když se dlouhá ukázka nevejde na slide, zmenši
% ji buď globálně tady, nebo jen v jedné prezentaci stejným řádkem za
% % ---------------------------------------------------------------------------
% Společná preambule pro prezentace předmětu Programování (3. ročník)
%
% Vzhled řeší motiv beamerthemeSPSProsek.sty ve stejné složce (převzatý
% z šablony předmětů ZPG/SNT na FS ČVUT a přeznačkovaný na SPŠ na Proseku).
% Tady se jen načte a doplní se metadata a makra předmětu.
%
% Překlad: XeLaTeX (kvůli fontspec a českým znakům)
%   latexmk -xelatex prezentace.tex
% nebo jednoduše  .\nastroje\build_prezentace.ps1
%
% Vkládá se z prezentace v temata/<tema>/prezentace/ takto:
%   \input{../../../spolecne/styl/preambule}
% ---------------------------------------------------------------------------

\usepackage{iftex}

% --- písma a čeština -------------------------------------------------------
\ifXeTeX
  \usepackage{fontspec}
  \usepackage{polyglossia}
  \setmainlanguage{czech}
  \setmainfont{Latin Modern Roman}
  \setsansfont{Latin Modern Sans}
  \setmonofont{Latin Modern Mono}[Scale=MatchLowercase]
\else
  % Záložní varianta, kdyby někdo přeložil pdfLaTeXem.
  \usepackage[T1]{fontenc}
  \usepackage[utf8]{inputenc}
  \usepackage{lmodern}
  \usepackage[czech]{babel}
\fi

\usepackage{graphicx}
\usepackage{booktabs}
\usepackage{xcolor}
\usepackage{tikz}
\usetikzlibrary{arrows.meta, positioning, shapes.geometric}

% --- motiv SPŠ na Proseku --------------------------------------------------
% Cesta ke složce se stylem a logy. Počítá s tím, že prezentace leží
% v temata/<tema>/prezentace/; zkusí se ale i mělčí zanoření, aby šlo
% přeložit i rozpracovaný soubor jinde v repozitáři.
\IfFileExists{../../../spolecne/styl/beamerthemeSPSProsek.sty}
  {\newcommand{\SPSStylPath}{../../../spolecne/styl}}
  {\IfFileExists{../../spolecne/styl/beamerthemeSPSProsek.sty}
    {\newcommand{\SPSStylPath}{../../spolecne/styl}}
    {\IfFileExists{../spolecne/styl/beamerthemeSPSProsek.sty}
      {\newcommand{\SPSStylPath}{../spolecne/styl}}
      {\newcommand{\SPSStylPath}{.}}}}
\usepackage{\SPSStylPath/beamerthemeSPSProsek}

\setbeamertemplate{footline}[SPSProsek]
\beamertemplatenavigationsymbolsempty

% --- velikost písma ve výpisech kódu ---------------------------------------
% Výchozí je \footnotesize. Když se dlouhá ukázka nevejde na slide, zmenši
% ji buď globálně tady, nebo jen v jedné prezentaci stejným řádkem za
% \input{...preambule}:
% \renewcommand{\SPSCodeSize}{\scriptsize}

% --- zkratky pro vkládání kódu ---------------------------------------------
\newcommand{\kod}[1]{\texttt{\small #1}}          % kód uvnitř věty
\newcommand{\kodSoubor}[1]{\lstinputlisting{#1}}  % celá skica ze souboru
\newcommand{\kodSouborRadky}[3]{\lstinputlisting[firstline=#2,lastline=#3]{#1}}

% --- upozornění na slidu ---------------------------------------------------
\newcommand{\pozor}[1]{\textcolor{akcent}{\textbf{#1}}}

% --- metadata předmětu -----------------------------------------------------
\author{Jaroslav Bušek}
\institute{Střední průmyslová škola na Proseku}
% \date{} zůstává prázdné záměrně -- datum a čas posledního překladu se
% zobrazuje automaticky v patě každého slidu (\SPSVersionStamp), takže je
% při promítání hned vidět, jak stará je verze PDF.
\date{}
:
% \renewcommand{\SPSCodeSize}{\scriptsize}

% --- zkratky pro vkládání kódu ---------------------------------------------
\newcommand{\kod}[1]{\texttt{\small #1}}          % kód uvnitř věty
\newcommand{\kodSoubor}[1]{\lstinputlisting{#1}}  % celá skica ze souboru
\newcommand{\kodSouborRadky}[3]{\lstinputlisting[firstline=#2,lastline=#3]{#1}}

% --- upozornění na slidu ---------------------------------------------------
\newcommand{\pozor}[1]{\textcolor{akcent}{\textbf{#1}}}

% --- metadata předmětu -----------------------------------------------------
\author{Jaroslav Bušek}
\institute{Střední průmyslová škola na Proseku}
% \date{} zůstává prázdné záměrně -- datum a čas posledního překladu se
% zobrazuje automaticky v patě každého slidu (\SPSVersionStamp), takže je
% při promítání hned vidět, jak stará je verze PDF.
\date{}
:
% \renewcommand{\SPSCodeSize}{\scriptsize}

% --- zkratky pro vkládání kódu ---------------------------------------------
\newcommand{\kod}[1]{\texttt{\small #1}}          % kód uvnitř věty
\newcommand{\kodSoubor}[1]{\lstinputlisting{#1}}  % celá skica ze souboru
\newcommand{\kodSouborRadky}[3]{\lstinputlisting[firstline=#2,lastline=#3]{#1}}

% --- upozornění na slidu ---------------------------------------------------
\newcommand{\pozor}[1]{\textcolor{akcent}{\textbf{#1}}}

% --- metadata předmětu -----------------------------------------------------
\author{Jaroslav Bušek}
\institute{Střední průmyslová škola na Proseku}
% \date{} zůstává prázdné záměrně -- datum a čas posledního překladu se
% zobrazuje automaticky v patě každého slidu (\SPSVersionStamp), takže je
% při promítání hned vidět, jak stará je verze PDF.
\date{}
:
% \renewcommand{\SPSCodeSize}{\scriptsize}

% --- zkratky pro vkládání kódu ---------------------------------------------
\newcommand{\kod}[1]{\texttt{\small #1}}          % kód uvnitř věty
\newcommand{\kodSoubor}[1]{\lstinputlisting{#1}}  % celá skica ze souboru
\newcommand{\kodSouborRadky}[3]{\lstinputlisting[firstline=#2,lastline=#3]{#1}}

% --- upozornění na slidu ---------------------------------------------------
\newcommand{\pozor}[1]{\textcolor{akcent}{\textbf{#1}}}

% --- metadata předmětu -----------------------------------------------------
\author{Jaroslav Bušek}
\institute{Střední průmyslová škola na Proseku}
% \date{} zůstává prázdné záměrně -- datum a čas posledního překladu se
% zobrazuje automaticky v patě každého slidu (\SPSVersionStamp), takže je
% při promítání hned vidět, jak stará je verze PDF.
\date{}
